\section{单个有限系统：从精确商到复杂度--误差前沿}
\label{sec:closure}

第~\ref{sec:operational}节固定了可辨微观空间 $B$，从而消除了仅由编码
复制产生的虚假压缩。本节先证明可辨约化严格保持原系统的闭合误差，再依次回答两个单系统问题：零误差时最少需要多少宏观状态；允许误差
以后，状态数与闭合误差之间的完整最优折中是什么。

对具有同一有限输入空间的 kernels 定义最坏行总变差距离
\[
d_\infty(K,K'):=\max_x\|K(x,\cdot)-K'(x,\cdot)\|_{\TV},
\qquad
\|\mu-\nu\|_{\TV}:=\frac12\sum_z|\mu(z)-\nu(z)|.
\]

\subsection{可辨约化保持闭合误差}

沿用命题~\ref{prop:operational-baseline} 的满射
$p:X\twoheadrightarrow B$。给定满射 $s:B\twoheadrightarrow Y$，令
$q=s\circ p$；$Q_p,Q_s,Q_q$ 表示相应的确定性 kernels。

\begin{proposition}[约化前后闭合缺陷相同]
\label{prop:operational-defect-invariance}
若 $KQ_p=Q_p\bar K$，则对每个 $L\in\Stoch(Y,Y)$，
\begin{equation}
 d_\infty(KQ_q,Q_qL)
 =d_\infty(\bar KQ_s,Q_sL).
 \label{eq:operational-defect-equality}
\end{equation}
因此，对 $L$ 取最小值、再对允许动力学取上确界时，等式仍成立。
\end{proposition}

\begin{proof}
由 $Q_q=Q_pQ_s$ 与 $KQ_p=Q_p\bar K$，等式两边所比较的两对 kernels
分别为 $Q_p\bar KQ_s$ 与 $Q_pQ_sL$，以及 $\bar KQ_s$ 与 $Q_sL$。
满射 $p$ 使前一对的行恰为后一对全部行的重复，故最大行距离相同。
\end{proof}

可辨约化同时规范复杂度分母并保持闭合缺陷：对所有按定义经 $B$ 因子化的候选
宏观表示，它严格保持后续的最优闭合误差。完整的代数构造与等式证明见
附录~\ref{app:operational-baseline}。

\subsection{精确闭合与唯一最粗任务商}

给定满射 $s:B\twoheadrightarrow Y$，若对
$\bar K\in\Stoch(B,B)$ 存在 $L\in\Stoch(Y,Y)$ 使
\begin{equation}
 \bar KQ_s=Q_sL,
 \label{eq:exact-closure}
\end{equation}
则称 $s$ 对 $\bar K$\emph{精确闭合}。这是对所有微观初始分布成立的
strong lumpability 条件\citep{KemenySnell1976,Buchholz1994}：同一
$s$-纤维中的状态必须对每个下一宏观块给出相同转移概率。

单个有限链的最粗 ordinary lumping 可由稳定分区细化构造
\citep{DerisaviEtAl2003}。以下有限精确结论同时预先保留任务，并要求同一个
表示服务于整个允许动力学族；每个微观 kernel 可以诱导自己的宏观 kernel。
该结论给出近似理论的零误差基准。
对函数 $f:B\to\mathbb R$，记
\[
(P_{\bar K}f)(b):=\sum_{b'\in B}\bar K(b,b')f(b')
\]
为 $\bar K$ 的观测算子。

\begin{theorem}[核族的唯一最粗精确任务商]
\label{thm:minimal-exact-task-quotient}
设 $B$ 为非空有限集，$\bar g:B\twoheadrightarrow O$ 为非恒定任务，
$\bar\cK\subseteq\Stoch(B,B)$ 非空。则存在一个在宏观状态重标记意义下
唯一的满射
\[
 s_*:B\twoheadrightarrow Y_*,
\]
满足：
\begin{enumerate}
  \item 存在唯一 $r_*:Y_*\to O$ 使 $\bar g=r_*\circ s_*$；
  \item 对每个 $\bar K\in\bar\cK$，存在唯一
  $L^*_{\bar K}\in\Stoch(Y_*,Y_*)$ 使
  \[
  \bar KQ_{s_*}=Q_{s_*}L^*_{\bar K};
  \]
  \item 若另一个任务保真满射 $s:B\twoheadrightarrow Y$ 对全部
  $\bar K\in\bar\cK$ 精确闭合，则存在唯一满射
  $h:Y\twoheadrightarrow Y_*$ 使 $s_*=h\circ s$。
\end{enumerate}
换言之，$s_*$ 是全部共同精确闭合、任务保真商中的唯一最粗者。定义
\begin{equation}
 c_{\bar\cK}^{0}(\bar g):=|Y_*|.
 \label{eq:exact-task-complexity}
\end{equation}
则严格精确宏观压缩存在，当且仅当
\begin{equation}
 |O|\le c_{\bar\cK}^{0}(\bar g)<|B|.
 \label{eq:strict-exact-compression}
\end{equation}
\end{theorem}

\begin{proof}
从任务观测代数 $\bar g^*\mathbb R^O$ 出发，反复加入所有
$P_{\bar K}$ 的像并取生成代数。有限维性保证迭代终止；稳定代数给出共同
精确闭合商。任何其他候选的分区代数都包含初始代数并对全部
$P_{\bar K}$ 不变，因而必包含该最小稳定代数，得到最粗性与唯一性。完整
普适性质和分区细化证明见附录~\ref{app:exact-theory}。
\end{proof}

一步交换式进一步给出全部整数时刻的半群交换。对每个固定
$\bar K\in\bar\cK$ 和 $t\in\mathbb N_0$，
\[
 \bar K^tQ_{s_*}=Q_{s_*}(L^*_{\bar K})^t,
\]
并且从任意微观初始分布出发，投影过程 $s_*(X_t)$ 本身是以
$L^*_{\bar K}$ 演化的 Markov 链。定理~\ref{thm:minimal-exact-task-quotient}
确定了复杂度--误差图的零误差端点。严格压缩恰在
$c_{\bar\cK}^{0}(\bar g)<|B|$ 时出现；恒等尺度
$c_{\bar\cK}^{0}(\bar g)=|B|$ 也可能成为该端点。

\subsection{允许误差后的复杂度--误差前沿}

固定状态预算下优化聚合误差是 Markov 状态约化中的一条相关路线；例如
\citet{GeigerEtAl2015}研究平稳 KL divergence rate 目标。本文采用任务保真
约束和最坏行 TV 缺陷，并把各预算下的最优值组织为如下前沿。

令 $\pi$ 为 $B$ 的一个任务保真分区，即每个 $\pi$-块都包含在某个
$\bar g$-纤维中；记商映射为
$s_\pi:B\twoheadrightarrow B/\pi$。定义该分区对整个动力学族的闭合缺陷
\begin{equation}
\eta_\pi(\bar\cK)
:=
\sup_{\bar K\in\bar\cK}
\min_{L\in\Stoch(B/\pi,B/\pi)}
d_\infty(\bar KQ_{s_\pi},Q_{s_\pi}L).
\label{eq:eta-pi}
\end{equation}
同一个 $\pi$ 服务于全部 $\bar K$，每个 $\bar K$ 拥有自己的最优 $L$；因此本文
使用 $\sup_{\bar K}\min_L$。整个动力学族共享同一个 $L$ 对应更强的
$\min_L\sup_{\bar K}$ 问题。命题~\ref{prop:operational-defect-invariance}
保证式~\eqref{eq:eta-pi} 与原始空间上复合表示 $s_\pi\circ p$ 的缺陷完全
相同。

对 $|O|\le m\le|B|$，令
\[
 \mathfrak P_m(\bar g)
 :=\{\pi:\pi\text{ 细化 }\ker\bar g,\ |\pi|\le m\},
\]
并定义
\begin{equation}
 \mathcal E_{\bar\cK,\bar g}(m)
 :=\min_{\pi\in\mathfrak P_m(\bar g)}\eta_\pi(\bar\cK).
 \label{eq:complexity-error-frontier}
\end{equation}
对容限 $\tau\ge0$，相应的任务相对最小状态数为
\begin{equation}
c_{\bar\cK}^{\tau}(\bar g)
:=\min\{ |\pi|:\pi\text{ 细化 }\ker\bar g,
                    \ \eta_\pi(\bar\cK)\le\tau\}.
\label{eq:relative-complexity}
\end{equation}
当核族为单点集 $\{\bar K\}$ 时，下标简记为 $\bar K$。

\begin{theorem}[复杂度--误差前沿的有限达到]
\label{thm:complexity-error-frontier}
对每个 $m$，式~\eqref{eq:complexity-error-frontier} 的最小值由某个任务
保真分区达到；对该分区和每个固定 $\bar K$，式~\eqref{eq:eta-pi} 内的
最优宏观 kernel 也达到。此外：
\begin{enumerate}
  \item $m\mapsto\mathcal E_{\bar\cK,\bar g}(m)$ 单调不增，且
  $\mathcal E_{\bar\cK,\bar g}(|B|)=0$；
  \item
  \[
  \mathcal E_{\bar\cK,\bar g}(m)=0
  \quad\Longleftrightarrow\quad
  c_{\bar\cK}^{0}(\bar g)\le m;
  \]
  \item 对任意 $\tau\ge0$，
  \begin{equation}
  c_{\bar\cK}^{\tau}(\bar g)
  =\min\{m:\mathcal E_{\bar\cK,\bar g}(m)\le\tau\}.
  \label{eq:tolerance-frontier-inverse}
  \end{equation}
\end{enumerate}
\end{theorem}

\begin{proof}
$B$ 只有有限个任务保真分区，固定分区的宏观 kernel 空间是紧单纯形，故
所述最小值均达到。增大 $m$ 只扩张可行分区集；离散分区在 $m=|B|$ 时
给出零误差。误差为零等价于对全部 $\bar K$ 精确闭合，再用
定理~\ref{thm:minimal-exact-task-quotient} 得到第二项。第三项由两边定义
直接得到。细节见附录~\ref{app:approximate-theory}。
\end{proof}

精确情形有唯一最粗商；当 $\tau>0$ 时，达到同一最优值的分区一般不可比较，
也通常不存在唯一最粗者。因此近似理论的规范对象是前沿
$\mathcal E_{\bar\cK,\bar g}$ 及每一点的最优解集合。

\subsection{纤维预测直径：消去未知宏观核}

对 $b\in B$ 定义其下一宏观块分布
\[
p^{\pi}_{\bar K,b}(C):=\bar K(b,C),
\qquad C\in\pi,
\]
并定义\emph{统一纤维预测直径}
\begin{equation}
D_{\bar\cK}(\pi)
:=
\sup_{\bar K\in\bar\cK}
\max_{s_\pi(b)=s_\pi(b')}
\bigl\|p^{\pi}_{\bar K,b}-p^{\pi}_{\bar K,b'}\bigr\|_{\TV}.
\label{eq:fiber-diameter}
\end{equation}
它完全由候选分区上的下一块分布构成，直接度量同一纤维内状态的预测分歧。

\begin{lemma}[纤维半径--直径夹逼]
\label{lem:radius-diameter}
对任意有限 $B$、任务保真分区 $\pi$ 和非空动力学族 $\bar\cK$，
\begin{equation}
\frac12D_{\bar\cK}(\pi)
\le
\eta_\pi(\bar\cK)
\le
D_{\bar\cK}(\pi).
\label{eq:radius-diameter}
\end{equation}
\end{lemma}

\begin{proof}
固定 $\bar K$ 与一个 $s_\pi$-纤维。宏观核对应行的最优选择，是该纤维内
下一块分布集合在总变差距离下的 Chebyshev 中心。任意包围半径至少为直径
的一半；任选集合内一点作中心，包围半径又不超过直径。逐纤维取最大、再对
$\bar K$ 取上确界即得结论。完整中心公式见附录~\ref{app:approximate-theory}。
\end{proof}

定义不含宏观 kernel 的直径前沿
\begin{equation}
 \mathcal D_{\bar\cK,\bar g}(m)
 :=\min_{\pi\in\mathfrak P_m(\bar g)}D_{\bar\cK}(\pi).
 \label{eq:diameter-frontier}
\end{equation}
该最小值同样因候选分区有限而达到。逐分区应用
引理~\ref{lem:radius-diameter} 后取最小值得
\begin{equation}
 \frac12\mathcal D_{\bar\cK,\bar g}(m)
 \le\mathcal E_{\bar\cK,\bar g}(m)
 \le\mathcal D_{\bar\cK,\bar g}(m).
 \label{eq:frontier-radius-diameter}
\end{equation}
所以两条前沿具有相同的零点与渐近趋零行为。单个有限系统总有一条达到的
前沿；下一节研究的是，这条前沿能否随系统规模增长而逼近低相对复杂度、低
闭合误差的原点。
